% Projeto adaptado do template original em Template_UFFSTex/.
% Dados básicos transcritos de tcc.conf; atualizar sob solicitação.
\documentclass[portuguese,oneside]{UFFStex}

\usepackage{graphicx}
\usepackage[hidelinks]{hyperref}
\usepackage[alf]{abntex2cite}

\author{Matheus Henrique Rodrigues da Costa}
% Título em inglês pendente para o final do TCC I.
\title{Avaliação do TabPFN para classificação litológica e predição da composição mineral a partir de dados espectrais de satélite}{}
\tipotrabalho{\tci}
\curso{\cc}
\orientador{Samuel da Silva Feitosa}
\coorientador{Jean Marcel de Almeida Espinoza}
% Palavras-chave em português e inglês pendentes para o final do TCC I.
\palavraschave{}{}
% A classe define UFFS, Campus Chapecó e Chapecó e usa o ano da compilação.
% O ano de entrega ainda precisa ser confirmado em tcc.conf.

\begin{document}

% Resumo e abstract serão inseridos aqui quando escritos e revisados.
% Os ambientes obrigatórios do template são resumo e abstract.
% Ficha catalográfica e folha de aprovação são exigidas para TCC II.
% Versão parcial com introdução e fundamentação teórica/revisão bibliográfica.
\tableofcontents

% Fonte principal: tcc/markdown/01-introducao.md.
\chapter{INTRODUÇÃO}
\label{chap:introducao}

\section{APRESENTAÇÃO}

Nos últimos anos, estudos têm aplicado \textit{machine learning} a dados de sensoriamento remoto para apoiar o mapeamento geológico e a estimativa de informações geoquímicas. Na Província Mineral de Carajás, por exemplo, métodos de \textit{machine learning} foram empregados no mapeamento litológico com dados geofísicos, de relevo e de sensoriamento remoto \cite{costa2019}. Outro estudo combinou dados Landsat-8 e Sentinel-2 com análises geoquímicas para estimar a concentração de ouro em uma instalação de rejeitos \cite{zhang2023}. O sensoriamento remoto permite observar grandes áreas, inclusive locais de difícil acesso, e pode ajudar a selecionar pontos para investigação posterior. As previsões obtidas, contudo, precisam ser verificadas com dados geológicos e análises físicas.

As informações espectrais registradas por sensores remotos dependem da interação entre a radiação eletromagnética e os materiais presentes na superfície terrestre. Características de absorção, como posição e profundidade, podem estar relacionadas à composição dos materiais; estudos também mostram relações entre dados espectrais e variáveis geoquímicas, embora essa relação dependa do sensor, da superfície observada e dos dados de referência \cite{bai2023,zhang2023}.

O TabPFN é um \textit{foundation model} para dados tabulares, pré-treinado com conjuntos de dados sintéticos. Durante a inferência, recebe os exemplos de treinamento como contexto para fazer previsões sobre novos exemplos, podendo ser aplicado a tarefas de classificação e regressão sem que seus parâmetros precisem ser treinados do zero para cada conjunto de dados \cite{hollmann2025}. Aplicações recentes já investigam o TabPFN em problemas próximos, como a previsão de propriedades do solo a partir de dados espectrais e a caracterização geotécnica \cite{barkov2026,saito2026}. Esses estudos mostram que o uso do modelo em dados relacionados ao solo e à geociência já começou. Ainda assim, permanece relevante investigar seu comportamento na classificação litológica e na predição de composição mineral usando dados espectrais obtidos por satélite, objetivos específicos deste trabalho.

Este trabalho propõe-se a avaliar o desempenho do TabPFN na classificação litológica e na predição da composição mineral a partir de dados espectrais de satélite. Para preparar os dados, pretende-se integrar e harmonizar conjuntos geológicos e geoquímicos de diferentes regiões do mundo, associando coordenadas, classes litológicas e medidas laboratoriais compatíveis aos valores de reflectância de imagens Sentinel-2. Entre as fontes previstas estão dados de áreas brasileiras mapeadas e em exploração, ainda dependentes de recebimento e inspeção. A seleção das imagens e das amostras deverá considerar sua correspondência espacial e temporal e as condições da superfície observada.

A definição dos alvos dependerá da disponibilidade de referências adequadas: concentrações de elementos ou óxidos serão tratadas como medidas geoquímicas, sem serem interpretadas automaticamente como proporções de minerais. Também será verificado se os rótulos litológicos provêm de observações de campo ou de mapas derivados de imagens, distinguindo a validação com referências independentes da reprodução de uma classificação anterior. A avaliação pretende considerar a distribuição espacial dos dados e incluir regiões não utilizadas no ajuste, conforme a cobertura dos conjuntos selecionados. Busca-se verificar o potencial de apoio à triagem inicial de áreas para investigação e disponibilizar no Kaggle os conjuntos derivados cuja redistribuição seja permitida, preservando a procedência e os procedimentos de preparação.

\section{PROBLEMA DE PESQUISA}

Em que medida o TabPFN pode classificar litologias e predizer a composição mineral a partir de dados espectrais de satélite associados a referências geológicas e laboratoriais, e como a integração de conjuntos harmonizados de diferentes regiões afeta seu desempenho em áreas não utilizadas no ajuste?

\section{HIPÓTESE DE PESQUISA}

A hipótese principal deste trabalho é que os dados espectrais de satélite, associados a referências adequadas, permitem ao TabPFN produzir previsões de litologia e composição mineral com capacidade preditiva mensurável em amostras não utilizadas no ajuste do modelo. Uma hipótese adicional é que a integração de dados harmonizados de diferentes regiões pode ampliar a representatividade do treinamento e melhorar o desempenho em uma região externa, em comparação com o uso de dados de uma única região. Esse possível ganho será objeto de avaliação, sem pressupor que a agregação de bases ou a diversidade de instrumentos necessariamente melhora as previsões. O potencial de apoio à priorização de áreas para investigação será discutido considerando os erros, a incerteza e as limitações observadas.

\section{OBJETIVOS}

\subsection{Objetivos gerais}

Avaliar o desempenho e a generalização do TabPFN na classificação litológica e na predição da composição mineral a partir de dados espectrais de satélite associados a referências geológicas e laboratoriais de diferentes regiões.

\subsection{Objetivos específicos}

\begin{itemize}
    \item Selecionar e harmonizar conjuntos geológicos e geoquímicos de diferentes regiões, verificando procedência, materiais, métodos, unidades, coordenadas e referências disponíveis para cada tarefa;
    \item Construir conjuntos tabulares pela integração espacial das referências selecionadas com valores de reflectância de imagens Sentinel-2, considerando compatibilidade temporal, resolução espacial e condições da superfície;
    \item Aplicar o TabPFN à classificação litológica e à predição da composição mineral nos conjuntos que possuam referências adequadas, documentando as limitações dos alvos disponíveis;
    \item Avaliar as previsões com métricas apropriadas e validação espacial e entre regiões, comparando, quando houver dados compatíveis suficientes, o treinamento com uma e com múltiplas regiões em uma mesma região externa de teste;
    \item Disponibilizar no Kaggle os conjuntos derivados cuja redistribuição seja autorizada pelas fontes, acompanhados de documentação sobre procedência, transformações e limitações;
    \item Analisar o potencial e as limitações das previsões para apoiar a priorização de áreas para investigação física.
\end{itemize}

\section{JUSTIFICATIVA}

O Brasil detém a segunda maior reserva de terras raras do mundo, mas responde por apenas 1\% da produção global; além disso, apenas 30\% do território brasileiro está mapeado \cite{g12026}. Esse cenário tem levado o poder público a atuar diretamente sobre a lacuna: em setembro de 2026, o governo federal sancionou uma política nacional voltada à ampliação da pesquisa, extração e processamento de minerais críticos e estratégicos, prevendo até R\$ 7 bilhões em instrumentos de incentivo e um aumento expressivo --- de R\$ 8 milhões para R\$ 300 milhões --- no orçamento destinado à investigação do subsolo brasileiro \cite{horaextra2026}. Nesse contexto, o desenvolvimento de métodos mais eficientes para a classificação litológica e a predição de composição mineral a partir de dados de sensoriamento remoto se torna especialmente relevante, uma vez que pode contribuir para acelerar o mapeamento geológico do país e subsidiar decisões estratégicas relacionadas à exploração desses recursos.

O TabPFN será investigado por ser um modelo voltado para dados tabulares, como os formados por bandas espectrais e índices de sensoriamento remoto. Assim, o trabalho busca verificar se essa abordagem pode apoiar o mapeamento inicial de áreas de interesse, sem substituir a confirmação geológica e ambiental em campo \cite{hollmann2025}.

A integração de bases de diferentes regiões também permite investigar a transferência do modelo para condições não representadas no ajuste. Barkov et al. (2026) reconhecem que coleções obtidas com métodos e instrumentos distintos introduzem variabilidade e que o desempenho em regiões ou instrumentos diferentes exige avaliação adicional; seus resultados não demonstram que a diversidade de sensores, por si só, melhora a predição \cite{barkov2026}. Neste trabalho, a harmonização dos registros e o uso pretendido de uma fonte orbital comum buscam construir conjuntos comparáveis para testar a hipótese de generalização. A documentação e a disponibilização dos derivados permitidos podem contribuir para a reprodução dos experimentos, independentemente de ser observado ganho de desempenho.

Do ponto de vista legal e institucional, o mapeamento de informações geológicas e minerais está relacionado às atribuições da Agência Nacional de Mineração (ANM), criada pela Lei nº 13.575/2017 para promover a gestão dos recursos minerais da União, além de regular e fiscalizar as atividades de aproveitamento desses recursos \cite{brasil2017}. O Plano Nacional de Mineração 2050 também reconhece a necessidade de ampliar o conhecimento geológico do território brasileiro e destaca a importância dos minerais críticos e estratégicos para a transição energética, a segurança alimentar, a tecnologia e a soberania nacional \cite{pnm2050}. Portanto, uma ferramenta de triagem baseada em dados espectrais pode contribuir para organizar etapas iniciais de investigação, embora não substitua os procedimentos legais, ambientais e técnicos exigidos para a pesquisa e a lavra mineral.

Em termos práticos, a prospecção tradicional depende de levantamentos de campo, análises laboratoriais e deslocamento de equipes, atividades que podem ser demoradas e custosas quando aplicadas a grandes áreas. O uso de sensoriamento remoto associado a modelos preditivos permite analisar uma quantidade maior de pontos e indicar onde a coleta física pode ser mais informativa. Dessa forma, o resultado esperado não é confirmar a existência de uma jazida, mas apoiar a seleção de áreas para investigação posterior.

\chapter{FUNDAMENTAÇÃO TEÓRICA E REVISÃO BIBLIOGRÁFICA}
\label{chap:fundamentacao}


\bibliography{bibliografia}
\end{document}
